% Generated draft: requires mathematical and source review.
\documentclass[11pt]{article}
\usepackage[margin=25mm]{geometry}
\usepackage{fontspec}
\setmainfont{Latin Modern Roman}
\usepackage{amsmath,amssymb,mathtools,unicode-math}
\setmathfont{Latin Modern Math}
\usepackage{enumitem,needspace,xurl,hyperref}
\setlength{\parindent}{0pt}
\setlength{\parskip}{5pt}
\setlength{\emergencystretch}{3em}
\newcommand{\N}{\mathbb{N}}
\newcommand{\Z}{\mathbb{Z}}
\newcommand{\Q}{\mathbb{Q}}
\newcommand{\R}{\mathbb{R}}
\newcommand{\C}{\mathbb{C}}
\newcommand{\F}{\mathbb{F}}
\newcommand{\A}{\mathbb{A}}
\newcommand{\PP}{\mathbb{P}}
\newcommand{\E}{\mathbb{E}}
\newcommand{\eps}{\varepsilon}
\newcommand{\dd}{\,\mathrm{d}}
\begin{document}
\section*{UnsolvedMath ID: 700228}
\textbf{Source problem code:} catalogue entry
\subsection{Definitions and mathematical statement}
Let \(\mathfrak K\) be a class of groups, understood up to isomorphism. For a group \(G\), let \(\operatorname{Aut}(G)\) denote its group of all automorphisms under composition. An epimorphic image means the target of a surjective group homomorphism. A countable group means a group whose underlying set has cardinality at most \(\aleph_0\).

Characterize all classes \(\mathfrak K\) satisfying the following conditions:
\begin{enumerate}
\item For every \(G\in\mathfrak K\) and every subgroup \(H\leq G\), one has \(H\in\mathfrak K\).
\item For every \(G\in\mathfrak K\), every group \(Q\), and every surjective homomorphism \(\varphi:G\twoheadrightarrow Q\), one has \(Q\in\mathfrak K\).
\item For every \(G\in\mathfrak K\), one has \(\operatorname{Aut}(G)\in\mathfrak K\).
\item There exists a group \(C\) such that \(|C|\leq\aleph_0\) and \(C\notin\mathfrak K\).
\end{enumerate}
The first and third conditions together also require every subgroup of \(\operatorname{Aut}(G)\) to belong to \(\mathfrak K\) whenever \(G\in\mathfrak K\), thus covering groups of automorphisms in that sense. The members of \(\mathfrak K\) are not themselves required to be countable. The source imposes no explicit nonemptiness condition on \(\mathfrak K\); under the literal conditions, the empty class is allowed.

The source states that the class of all finite groups and the class of all almost cyclic groups satisfy these requirements. The supplied record does not define ``almost cyclic,'' so the precise meaning of that second example remains unresolved; no definition is imposed here.
\subsection{Short English statement}
Characterize the classes of groups closed under taking subgroups, surjective homomorphic images, and automorphism groups, while omitting at least one countable group. The source gives finite groups and ``almost cyclic'' groups as examples, but does not define the latter term.
\subsection{Sources}
\begin{itemize}
\item Supplied local UnsolvedMath registry record, permanent ID \(700228\), exact source problem code \texttt{catalogue entry}, title \emph{Kourovka Notebook Problem 4.18}; upstream revision \texttt{d371ee917ac9647fbb06cd6fc529deec6f3dda82}. The supplied registry labels the problem ``open''; this label is not an independently verified mathematical or literature assessment.
\item R. Baer, Problem 4.18, \emph{The Kourovka Notebook (Unsolved Problems in Group Theory)}, fourth issue (1973), as attributed in the supplied record.
\item The supplied source identifies the twenty-first issue edition (2026), arXiv:1401.0300v48. The collection metadata names E. I. Khukhro and V. D. Mazurov as editors of the Issues 1--21 collection (1965--2026). These bibliographic details are supplied attribution and have not been independently checked.
\end{itemize}
\end{document}
